\subsection{依测度收敛}

设 X 是一个局部紧空间，\(\mu\) 是 X 上的一个测度，A 是一个 \(\mu\)-可测的、含于 X 的子集，F 是一个一致空间；我们将用 \(\mathcal{S}(\mathrm{A},\mu;\mathrm{F})\)，或 \(\mathcal{S}_{\mathrm{F}}(\mathrm{A},\mu)\)（或当 A = X 时简写为 \(\mathcal{S}_{\mathrm{F}}(\mu)\)，甚至 \(S_{F}\)）表示 A 到 F 中的 \(\mu\)-可测映射所成之集（No.~10, 定义 8）。对 F 的一致结构的每个近域 V、每个 \(\mu\)-可积集 \(B \subset A\) 以及每个数 \(\delta > 0\)，我们用 \(\boldsymbol{W}(\mathrm{V},\mathrm{B},\delta)\) 表示具有下列性质的函数对 \((f,g)\)（\(\mathcal{S}(\mathrm{A},\mu;\mathrm{F})\) 中的）所成之集：由所有 \(x \in B\)（使 \((f(x),g(x)) \notin V\) 者）所成之集 M 满足 \(|\mu|^{*}(\mathrm{M}) \leqslant \delta\) 。我们来证明诸集 \(\boldsymbol{W}(\mathrm{V},\mathrm{B},\delta)\) 构成 \(\mathcal{S}(\mathrm{A},\mu ;\mathrm{F})\) 上某个一致结构的基本近域系：显然，若 V 对称则 \(W(\mathrm{V},\mathrm{B},\delta)\) 对称，且若 \(\mathrm{V}'\subset \mathrm{V}\) 、 \(\mathrm{B}'\supset \mathrm{B}\) 与 \(\delta^{\prime}\leqslant \delta\)，则

\[
\boldsymbol {W} \left(\mathrm{V} ^ {\prime}, \mathrm{B} ^ {\prime}, \delta^ {\prime}\right) \subset \boldsymbol {W} (\mathrm{V}, \mathrm{B}, \delta);
\]

因此只需验证公理 \((\mathrm{U}_{\mathrm{III}}^{\prime})\)（GT, II, §1, No.~1）。而若 \(V'\) 是使 \(V^{\prime} \subset V\) 的一个近域，则

\[
\boldsymbol {W} \left(\mathrm{V} ^ {\prime}, \mathrm{B}, \delta / 2\right) \circ \boldsymbol {W} \left(\mathrm{V} ^ {\prime}, \mathrm{B}, \delta / 2\right) \subset \boldsymbol {W} (\mathrm{V}, \mathrm{B}, \delta).
\]

注意，当 K 取遍 A 的紧子集所成的一个 \(\mu\)-稠密集族 \(\mathfrak{K}\) 时，诸集 \(\boldsymbol{W}(\mathrm{V},\mathrm{K},\delta)\) 也构成上述一致结构的一个基本近域系：因为，对每个可积集 \(B\subset A\)，存在含于 B 的紧集 \(K\in \mathfrak{K}\)，使 \(|\mu|(\mathrm{B}-\mathrm{K})\leqslant\delta\)，从而 \(\boldsymbol{W}(\mathrm{V},\mathrm{K},\delta)\subset\boldsymbol{W}(\mathrm{V},\mathrm{B},2\delta)\) 。

定义 9. --- 在 \(\mathcal{S}(\mathrm{A},\mu;\mathrm{F})\) 上以诸 \(\boldsymbol{W}(\mathrm{V},\mathrm{B},\delta)\) 为基本近域系的一致结构，称为 A 中依测度收敛的一致结构。

相应的拓扑称为 A 中依测度收敛的拓扑，而对该拓扑收敛的滤子（或序列）称为在 A 中依测度收敛；当 A = X 时常略去对 A 的提及。

设 F 是 Hausdorff 的；那么，对每个 \(\mu\)-可积集 \(B \subset A\)，诸近域 \(\boldsymbol{W}(V, B, \delta)\)（其中 V 取遍 F 的一个基本近域系，\(\delta\) 取遍 \textgreater0 的数所成之集）的交，是由这样的对 \((f, g)\)（使 \(f(x) = g(x)\) 在 B 中（关于 \(\mu\)）几乎处处成立者）组成的。因为，由 \(x \in B\)（满足 \(f(x) \neq g(x)\) 者）所成之集 M 是 \(\mu\)-可积的，因为它是 \(\mu\)-可测映射 \(x \mapsto (f(x), g(x))\) 在 \(F \times F\) 中对角集的（开）补集下的原像（No.~5, 命题 7）；若 \(|\mu|(M) = \alpha > 0\)，则存在紧子集 \(K \subset M\)，使 \(|\mu|(M - K) < \alpha/2\)，且 f 与 g 在 K 上的限制连续；因此存在 F 的一个近域 \(V_0\)，使 \((f(x), g(x)) \notin V_0\) 对所有 \(x \in K\) 成立，从而 \((f, g) \notin W(V_0, B, \alpha/2)\) 。

由此推出，若 F 是 Hausdorff 的，则 \(\mathcal{S}(\mathrm{A},\mu;\mathrm{F})\) 的所有近域之交是由这样的对 \((f,g)\)（使 \(f(x)=g(x)\) 在 A 中局部几乎处处成立者）组成的。于是，与 \(\mathcal{S}(\mathrm{A},\mu;\mathrm{F})\) 相伴的 Hausdorff 一致空间——我们将把它记作 \(\mathrm{S}(\mathrm{A},\mu;\mathrm{F})\) 或 \(\mathrm{S}_{\mathrm{F}}(\mathrm{A},\mu)\)（当 A=X 时甚至记作 \(\mathrm{S}_{\mathrm{F}}(\mu)\) 或 \(S_{F}\)）——由对关系式 \(\langle f(x)=g(x)\) 在 A 中局部几乎处处成立 \(\rangle\) 的等价类（在集合 \(\mathcal{S}(\mathrm{A},\mu;\mathrm{F})\) 中）组成。

命题 17. --- 设 \((\mathrm{A}_{\lambda})_{\lambda \in \mathrm{L}}\) 是 A 的 \(\mu\)-可测子集所成的一个局部可数族，诸集两两不相交，且使 \(\mathrm{A} - \bigcup_{\lambda \in \mathrm{L}} \mathrm{A}_{\lambda}\) 局部 \(\mu\)-可忽略。若对每个类 \(\dot{f} \in \mathrm{S}(\mathrm{A}, \mu; \mathrm{F})\) 与每个 \(\lambda \in \mathrm{L}\)，用 \(\dot{f}_{\lambda}\) 表示到 \(\mathrm{A}_{\lambda}\) 上的限制（类 \(\dot{f}\) 中任一函数的限制）所在的类，则映射 \(\psi: \dot{f} \mapsto (\dot{f}_{\lambda})_{\lambda \in \mathrm{L}}\) 是一致空间 \(\mathrm{S}(\mathrm{A}, \mu; \mathrm{F})\) 到积一致空间 \(\prod_{\lambda \in \mathrm{L}} \mathrm{S}(\mathrm{A}_{\lambda}, \mu; \mathrm{F})\) 上的一个同构。

由 No.~10 命题 16，\(\psi\) 是双射。考虑 S(A, \(\mu\) ; F) 的一个近域 T，它是某个 \(W(V, B, \delta)\)（B 为 A 的紧子集）的规范像；我们知道，由 \(\lambda \in L\) 中使 B \(\cap\) A \(_{\lambda}\) \(\neq\) \(\varnothing\) 者所成之集 J 是可数的（No.~9），且 \(|\mu|(B) = \sum_{\lambda \in J} |\mu|(B \cap A_{\lambda})\)；因此存在 J 的一个有限子集 H，使得

\[
\sum_ {\lambda \in \mathrm{J} - \mathrm{H}} | \mu | (\mathrm{B} \cap \mathrm{A} _ {\lambda}) \leqslant \frac {\delta}{2}.
\]

于是 T 在 \(\psi \times \psi\) 下的像含于积 \(\prod_{\lambda \in H} W(V, B \cap A_{\lambda}, \delta)\) 的规范像。另一方面，若 \(m\) 是 H 的元素个数，则 T 在 \(\psi \times \psi\) 下的像包含近域 \(\prod_{\lambda \in H} W(V, A_{\lambda}, \delta / 2m)\) 的规范像，这就证明了命题。

命题 18. --- 若 F 可度量化，且 A 是一个局部 \(\mu\)-可忽略集与一个序列 \((A_{n})\)（由 \(\mu\)-可积集组成）之并，则空间 \(S(A,\mu;F)\) 可度量化。

由于每个 \(A_{n}\) 是一个可忽略集与一个紧集序列之并，我们可以假设诸 \(A_{n}\) 已经是紧的且两两不相交。于是命题 17 使我们可以假设 A 是紧的。若 \((\mathrm{V}_{n})\) 是 F 的一个可数基本近域系，则当 n 取遍 N 时，诸 \(\boldsymbol{W}(\mathrm{V}_{n},\mathrm{A},1/n)\) 显然构成 \(\mathcal{S}(\mathrm{A},\mu;\mathrm{F})\) 的一个基本近域系，命题得证。

引理 4. --- 设 F 是一个可度量化的一致空间，B ⊂ A 是 μ-可积集的一个可数并。则对每个 Cauchy 序列 ( \(f_{n}\) )（\(\mathcal{S}(A,\mu;F)\) 中的），存在一个子序列 ( \(f_{n_{k}}\) )（( \(f_{n}\) ) 的），使 ( \(f_{n_{k}}(x)\) ) 对几乎每个 \(x \in B\) 是 F 中的 Cauchy 序列。

先设 B 可积，并用 d 表示与 F 的一致结构相容的一个距离。我们将递归地定义一个二重序列 \((f_{mn})\)（由 \(\mathcal{S}(\mathrm{A},\mu;\mathrm{F})\) 中的函数组成），使得对每个 n 有 \(f_{0n}=f_{n}\)，\((f_{mn})_{n\geqslant0}\) 是 \((f_{m-1,n})_{n\geqslant0}\) 的一个子序列（对每个 m\textgreater0），最后，使得对每个 m\textgreater0，集 \(M_{mn}\)——由 \(x\in B\)（满足 \(d(f_{mn}(x),f_{m,n+1}(x))>1/2^{m+n+1}\) 者）组成——的测度 \(|\mu|(\mathrm{M}_{mn})\leqslant1/2^{m+n+1}\)；

这样一个定义的可行性来自 \((f_n)\) 是 \(\mathcal{S}(\mathrm{A},\mu ;\mathrm{F})\) 中的 Cauchy 序列这一事实。令 \(\mathbf{M}_m = \bigcup_{n\geqslant 0}\mathbf{M}_{mn}\)；则

\[
| \mu | (\mathrm{M} _ {m}) \leqslant \sum_ {n = 0} ^ {\infty} | \mu | (\mathrm{M} _ {m n}) \leqslant 1 / 2 ^ {m}
\]

且对每个 \(x \in \mathrm{B} - \mathrm{M}_m\)，有 \(d(f_{mn}(x), f_{m,n+p}(x)) \leqslant 1/2^{m+n}\) 对所有 \(n \geqslant 0\) 与所有 \(p > 0\) 成立；于是序列 \((f_{mn}(x))_{n \geqslant 0}\) 是 F 中的 Cauchy 序列。现在令 \(\mathbf{N} = \bigcap_{m=0}^{\infty} \mathbf{M}_m\)；N 是可忽略的。令 \(g_n = f_{nn}\)（对每个 \(n \geqslant 0\)）；对每个 \(m\)，序列 \((g_n)_{n \geqslant m}\) 是序列 \((f_{mn})_{n \geqslant 0}\) 的一个子序列；若 \(x \in \mathrm{B} - \mathrm{N}\)，则存在指标 \(m\) 使 \(x \notin \mathrm{M}_m\)，这就证明序列 \((g_n(x))\) 是 F 中的 Cauchy 序列。

若现在 B 是可积集序列 \((\mathrm{B}_{m})\) 之并，则可以递归地定义一个二重序列 \((g_{mn})\)，使得 \(g_{0n} = f_{n}\)，\((g_{mn})_{n \geqslant 0}\) 是 \((g_{m-1,n})_{n \geqslant 0}\) 的一个子序列（对每个 m \textgreater{} 0），且序列 \((g_{mn}(x))_{n \geqslant 0}\) 在 \(B_{m} - P_{m}\) 中是 Cauchy 序列，其中 \(P_{m}\) 可忽略。令 \(h_{n} = g_{nn}\)（对每个 \(n \geqslant 0\)），使得对每个 m，序列 \((h_{n})_{n \geqslant m}\) 是 \((g_{mn})_{n \geqslant 0}\) 的一个子序列；于是序列 \((h_{n}(x))_{n \geqslant 0}\) 对每个 \(x \in B - P\) 是 F 中的 Cauchy 序列，其中 \(P = \bigcup_{m=0}^{\infty} P_{m}\) 可忽略。

命题 19. --- 若一致空间 F 可度量化且完备，则一致空间 S(A, μ; F) 完备。

存在 A 的紧子集所成的一个局部可数族 \((\mathrm{K}_{\lambda})_{\lambda \in L}\)，使诸 \(\mathbf{K}_{\lambda}\) 两两不相交，且 \(\mathbf{A} - \bigcup_{\lambda}\mathbf{K}_{\lambda}\) 局部可忽略（No.~9, 命题 14）。由命题 17，\(\mathrm{S}(A,\mu;F)\) 同构于积 \(\prod_{\lambda \in L}\mathrm{S}(K_{\lambda},\mu;F)\)；这样我们就归结为证明 A 可积时的命题；这时 \(\mathrm{S}(A,\mu;F)\) 可度量化（命题 18），且由引理 4，对每个 Cauchy 序列 \((f_n)\)（\(\mathcal{S}(\mathrm{A},\mu;\mathrm{F})\) 中的），存在一个子序列 \((f_{n_k})\)，它在 \(\mathbf{A} - \mathbf{N}\) 中收敛，其中 \(\mathbf{N}\) 可忽略；极限 \(f\)（\((f_{n_k})\) 的，以任意方式延拓到整个 A）于是是 \(\mu\)-可测的，且由 No.~10 中所述的 Egoroff 定理的推广可知，序列 \((f_{n_k})\) 在 A 中依测度收敛于 \(f\)。这蕴涵 \(f\) 是序列 \((f_n)\) 在 \(\mathcal{S}(\mathrm{A},\mu;\mathrm{F})\) 中的一个聚点，而由假设序列 \((f_n)\) 是 Cauchy 序列，故它收敛于 \(f\) 。

证毕。

推论. --- 设 F 是一个可度量化的一致空间。

\begin{enumerate}
\def\labelenumi{(\roman{enumi})}
\item
  每个序列 \((f_n)\)（由 \(\mathcal{S}(\mathrm{A},\mu ;\mathrm{F})\) 的元素组成），若它局部几乎处处收敛到 A 到 F 中的一个映射 \(f\)（它必然 \(\mu\)-可测），则在 A 中依测度收敛于 \(f\)。
\item
  设 \((f_n)\) 是 \(\mathcal{S}(\mathrm{A},\mu ;\mathrm{F})\) 中元素的一个序列，它依测度收敛到一个映射 \(f\)（\(\mathrm{A}\) 到 \(\mathrm{F}\) 中的）。对每个是可积集的可数并的集 \(\mathbf{B}\subset \mathbf{A}\)，存在一个子序列 \((f_{n_k})\)（\((f_n)\) 的），使序列 \((f_{n_k}(x))\) 在 \(\mathbf{F}\) 中收敛于 \(f(x)\)（对几乎每个 \(x\in \mathbf{B}\)）。
\setcounter{enumi}{0}
\item
  该断言由 No.~10 中所述的 Egoroff 定理的推广直接得出。
\item
  由引理 4，存在一个子序列 \((f_{n_k})\)（\((f_n)\) 的），使 \((f_{n_k}(x))\) 是 \(\mathbf{F}\) 中的 Cauchy 序列（对每个 \(x\in \mathrm{B} - \mathrm{N}\)），其中 \(\mathbf{N}\) 可忽略；设 \(f^{\prime}(x)\in \widehat{\mathrm{F}}\) 是该序列对 \(x\in \mathrm{B} - \mathrm{N}\) 的极限。显然 \(f^{\prime}\) 是一个 \(\mu\)-可测映射（\(\mathrm{B} - \mathrm{N}\) 到 \(\widehat{\mathrm{F}}\) 中的），且由 (i)，序列 \((f_n)\) 依测度收敛于 \(f^{\prime}\)（在 \(\mathrm{B} - \mathrm{N}\) 中）；因此 \(f^{\prime}\) 几乎处处等于 \(f\)（在 \(\mathrm{B}\) 中）。
\end{enumerate}

命题 20. --- 设 F 是一个 Banach 空间，赋以由其范数定义的一致结构。

\begin{enumerate}
\def\labelenumi{(\roman{enumi})}
\item
  对 X 的每个 \(\mu\)-可测子集 A，依测度收敛的拓扑与 \(\mathcal{S}(\mathrm{A},\mu;\mathrm{F})\) 的向量空间结构相容。
\item
  空间 \(\mathcal{K}(\mathrm{X};\mathrm{F})\) 在 \(\mathcal{S}(\mathrm{X},\mu ;\mathrm{F})\) 中稠密
\item
  对每个实数 \(p \geqslant 1\)，依测度收敛的拓扑在空间 \(\mathcal{L}_{\mathrm{F}}^{p}(\mathrm{X},\mu)\) 上诱导的拓扑粗于 \(p\) 阶平均收敛的拓扑。
\setcounter{enumi}{0}
\item
  对 A 的每个 \(\mu\)-可积子集 B 与每个 \(\delta > 0\)，用 \(T(B, \delta)\) 表示这样的 \(\mathbf{f} \in \mathcal{S}(A, \mu; F)\) 所成之集：由 \(x \in B\)（满足 \(|\mathbf{f}(x)| \geqslant \delta\) 者）所成之集 C 适合关系 \(|\mu|(C) \leqslant \delta\)；显然，若 \(V_{\delta}\) 是由对 \((\mathbf{y}, \mathbf{z})\)（满足 \(|\mathbf{y} - \mathbf{z}| \leqslant \delta\) 者）组成的 F 的近域，则近域 \(W(V_{\delta}, B, \delta)\) 是 A 到 F 中的可测映射的对 \((\mathbf{f}, \mathbf{g})\) 中使 \(\mathbf{f} - \mathbf{g} \in T(B, \delta)\) 者所成之集。显然诸集 \(T(B, \delta)\) 是对称的，且 \(T(B, \delta) + T(B, \delta) \subset T(B, 2\delta)\) 与 \(T(B, |\alpha| \delta) \subset \alpha T(B, \delta)\)（对每个非零标量 \(\alpha\)，满足 \(|\alpha| \leqslant 1\) 者）成立；因此只需验证诸集 \(T(B, \delta)\) 是吸收的（TVS, I, §1, No.~5, 命题 4）。而若 \(\mathbf{f}\) 是 A 到 F 中的一个 \(\mu\)-可测映射，则数值函数 \(|\mathbf{f}|\) 也是 \(\mu\)-可测的（No.~3, 定理 1 的推论 6）。设 \(C_n\) 是由 \(x \in B\)（满足 \(|\mathbf{f}(x)| \geqslant n\) 者）所成之集；诸 \(C_n\) 构成一个可积集的递减序列，其交为空；因此存在整数 \(n\)，使 \(|\mu|(C_n) \leqslant \delta\)（ \(\S 4\) , No.~5, 命题 7 的推论）；此外还可以假设 \(n\) 取得足够大，使 \(1/n \leqslant \delta\)；于是 \(f/n^2 \in T(B, \delta)\)，这就证明了断言 (i)。
\setcounter{enumi}{2}
\item
  关系式 \(\int |\mathbf{f}|^p d|\mu| \leqslant \delta^{p+1}\) 蕴涵：若 C 是由 \(x \in \mathbf{X}\)（满足 \(|\mathbf{f}(x)| \geqslant \delta\) 者）所成之集，则
\end{enumerate}

\[
\delta^ {p} | \mu | ^ {*} (\mathrm{C}) \leqslant \int | \mathbf {f} | ^ {p} d | \mu | \leqslant \delta^ {p + 1},
\]

由此得 \(|\mu|^*(\mathrm{C}) \leqslant \delta\)，这就证明了 (iii)。

\begin{enumerate}
\def\labelenumi{(\roman{enumi})}
\setcounter{enumi}{1}
\tightlist
\item
  鉴于 (iii)，例如只需证明 \(\mathcal{L}_{\mathrm{F}}^{1}\) 在 \(\mathcal{S}_{\mathrm{F}}\) 中稠密，因为按定义，\(\mathcal{K}(\mathrm{X};\mathrm{F})\) 在 \(\mathcal{L}_{\mathrm{F}}^{1}\) 中关于平均收敛的拓扑稠密。现在，设 \(\mathbf{f}\) 是 \(\mathcal{S}_{\mathrm{F}}\) 的任一元素，\(T(\mathrm{B},\delta)\) 是该空间中 0 的一个邻域；我们如同 (i) 那样看出，存在 B 的一个可积子集 C，使 \(|\mu|(\mathrm{C})\leqslant \delta\)，且 \(\mathbf{f}\) 在 \(\mathrm{B - C}\) 上有界；于是用 \(\mathbf{g}\) 表示等于 \(\mathbf{f}\)（在 \(\mathrm{B - C}\) 上）、等于 0（在 \(\mathrm{X - (B - C)}\) 上）的函数，则由 No.~6 定理 5 得 \(\mathbf{g}\) 可积，且显然 \(\mathbf{f} - \mathbf{g}\in T(\mathrm{B},\delta)\) 。
\end{enumerate}

注. --- 1) 拓扑向量空间 \(\mathcal{S}(\mathrm{X},\mu ;\mathrm{F})\) 未必是局部凸的（习题 24）。

\begin{enumerate}
\def\labelenumi{\arabic{enumi})}
\setcounter{enumi}{1}
\tightlist
\item
  依测度收敛的拓扑在使 \(\mathrm{N}_{p}(\mathbf{f}) \leqslant 1\) 的 f 所成之集上诱导的拓扑，可能严格粗于 p 阶平均收敛的拓扑在该集上诱导的拓扑（习题 22）。但参见下文的命题 21。
\end{enumerate}

定义 10. --- 设 X 是一个局部紧空间，\(\mu\) 是 X 上的一个测度，F 是一个 Banach 空间，\(p \in [1, +\infty[\) 。称 \(\mathcal{L}_{\mathrm{F}}^{p}(\mathrm{X}, \mu)\) 的一个子集 H 是 \(p\) 阶等度可积的（关于 \(\mu\) ），如果它满足下列条件：

\begin{enumerate}
\def\labelenumi{(\roman{enumi})}
\tightlist
\item
  对每个 \(\varepsilon > 0\)，存在一个 \(\delta > 0\)，使对每个测度 \(|\mu| (\mathrm{A}) \leqslant \delta\) 的可积集 A 与每个 \(\mathbf{f} \in \mathrm{H}\)，
\end{enumerate}

\[
\int | \mathbf {f} | ^ {p} \varphi_ {\mathrm{A}} d | \mu | \leqslant \varepsilon .
\]

\begin{enumerate}
\def\labelenumi{(\roman{enumi})}
\setcounter{enumi}{1}
\tightlist
\item
  对每个 \(\varepsilon > 0\)，存在一个紧子集 \(K\)（\(X\) 的），使对每个 \(\mathbf{f} \in H\)，\(\int |\mathbf{f}|^p \varphi_{X - K} d|\mu| \leqslant \varepsilon\) 。
\end{enumerate}

当 \(p = 1\) 时，我们说“等度可积”以代替“1 阶等度可积”。

注. --- 设 \(\mu\) 有界。对每个 a \textgreater{} 0，由使 \(|\mathbf{f}(x)| \leqslant a\) 几乎处处成立的 X 到 F 中的可测映射所成之集是 p 阶等度可积的，而且这对任何 \(p \in [1, +\infty[\) 都成立。

命题 21. --- 设 H 是 \(\mathcal{L}_{\mathrm{F}}^{p}(\mathrm{X},\mu)\) 的一个 \(p\) 阶等度可积的子集。在 H 上，依测度收敛的一致结构等于由 \(\mathcal{L}_{\mathrm{F}}^{p}(\mathrm{X},\mu)\) 的一致结构诱导的一致结构。

设 \(\varepsilon > 0\) 。存在具有定义 10 的性质 (i) 与 (ii) 的 \(\delta\) 与 \(K\)。设 \(\mathbf{f}, \mathbf{g}\)（\(H\) 中的）满足

\[
| \mathbf {f} (x) - \mathbf {g} (x) | \leqslant \left(\frac {\varepsilon}{| \mu | (\mathrm{K})}\right) ^ {1 / p}
\]

对 \(x \in \mathbf{K}\) 成立，但集 \(\mathbf{M}\)（测度 \(\leqslant \delta\) 的）上的点除外。则

\[
\left(\int_ {\mathrm{X} - \mathrm{K}} | \mathbf {f} - \mathbf {g} | ^ {p} d | \mu |\right) ^ {1 / p} \leqslant \left(\int_ {\mathrm{X} - \mathrm{K}} | \mathbf {f} | ^ {p} d | \mu |\right) ^ {1 / p} + \left(\int_ {\mathrm{X} - \mathrm{K}} | \mathbf {g} | ^ {p} d | \mu |\right) ^ {1 / p}
\]

且同理

\[
\left(\int_ {\mathrm{M}} | \mathbf {f} - \mathbf {g} | ^ {p} d | \mu |\right) ^ {1 / p} \leqslant 2 \varepsilon^ {1 / p},
\]

因此

\[
\begin{array}{c} \int | \mathbf {f} - \mathbf {g} | ^ {p} d | \mu | = \int_ {\mathrm{X-K}} | \mathbf {f} - \mathbf {g} | ^ {p} d | \mu | + \int_ {\mathrm{M}} | \mathbf {f} - \mathbf {g} | ^ {p} d | \mu | + \int_ {\mathrm{K-M}} | \mathbf {f} - \mathbf {g} | ^ {p} d | \mu | \\ \leqslant 2 ^ {p} \varepsilon + 2 ^ {p} \varepsilon + \frac {\varepsilon}{| \mu | (\mathrm{K})} | \mu | (\mathrm{K-M}) \leqslant (2 ^ {p + 1} + 1) \varepsilon . \end{array}
\]

于是，H 上依测度收敛的一致结构细于由 \(\mathcal{L}_{\mathrm{F}}^{p}(\mathrm{X},\mu)\) 的一致结构诱导的一致结构。接下来只需应用命题 20。

